% Part: normal-modal-logic
% Chapter: tableaux
% Section: rules-for-K

\documentclass[../../../include/open-logic-section]{subfiles}

\begin{document}

\olfileid{nml}{tab}{rul}

\olsection{\Ax{K} دے قاعدے}

عام قضیاتی رابطیاں لئی قاعدے اوہی نیں جیہڑے عام قضیاتی نشان والے
!!{tableau}s لئی نیں، بس اگیتے شامل ہو گئے نیں۔ ہر صورت وچ، نشان والے !!{formula}
$\sFmla{S}{!A}[\sigma]$ اُتے قاعدہ لاؤن نال نویں !!{formula}s ملدے نیں
جیہناں دا اگیتا وی~$\sigma$ اے۔ ایہ گل سِدھی اے۔ مثلاً جے $!A
\land !B$،~$\sigma$ دے ناں والے جہان وچ سچ اے، تاں $!A$ تے $!B$
اوسے~$\sigma$ وچ سچ نیں (کسی ہور جہان وچ نہیں)۔ قضیاتی قاعدے اسیں
\olref{tab:prop-rules} وچ اکٹھے دتے نیں۔

\begin{table}
  \[\def\arraystretch{3}\begin{array}{|c|c|}
    \hline
    \AxiomC{\sFmla{\True}{\lnot !A}[\sigma]}
    \RightLabel{\TRule{\True}{\lnot}}
    \UnaryInfC{\sFmla{\False}{!A}[\sigma]}
    \DisplayProof
    &
    \AxiomC{\sFmla{\False}{\lnot !A}[\sigma]}
    \RightLabel{\TRule{\False}{\lnot}}
    \UnaryInfC{\sFmla{\True}{!A}[\sigma]}
    \DisplayProof
    \\[1ex]
    \hline
    \AxiomC{\sFmla{\True}{!A \land !B}[\sigma]}
    \RightLabel{\TRule{\True}{\land}}
    \UnaryInfC{\sFmla{\True}{!A}[\sigma]}
    \noLine
    \UnaryInfC{\sFmla{\True}{!B}[\sigma]}
    \DisplayProof
    &
    \AxiomC{\sFmla{\False}{!A \land !B}[\sigma]}
    \RightLabel{\TRule{\False}{\land}}
    \UnaryInfC{$\sFmla{\False}{!A}[\sigma] \quad \mid \quad
      \sFmla{\False}{!B}[\sigma]$}
    \DisplayProof
    \\[2ex]
    \hline
    \AxiomC{\sFmla{\True}{!A \lor !B}[\sigma]}
    \RightLabel{\TRule{\True}{\lor}}
    \UnaryInfC{$\sFmla{\True}{!A}[\sigma] \quad \mid \quad
      \sFmla{\True}{!B}[\sigma]$}
    \DisplayProof
    &
    \AxiomC{\sFmla{\False}{!A \lor !B}[\sigma]}
    \RightLabel{\TRule{\False}{\lor}}
    \UnaryInfC{\sFmla{\False}{!A}[\sigma]}
    \noLine
    \UnaryInfC{\sFmla{\False}{!B}[\sigma]}
    \DisplayProof
    \\[2ex]
    \hline
    \AxiomC{\sFmla{\True}{!A \lif !B}[\sigma]}
    \RightLabel{\TRule{\True}{\lif}}
    \UnaryInfC{$\sFmla{\False}{!A}[\sigma] \quad \mid
      \quad \sFmla{\True}{!B}[\sigma]$}
    \DisplayProof
    &
    \AxiomC{\sFmla{\False}{!A \lif !B}[\sigma]}
    \RightLabel{\TRule{\False}{\lif}}
    \UnaryInfC{\sFmla{\True}{!A}[\sigma]}
    \noLine
    \UnaryInfC{\sFmla{\False}{!B}[\sigma]}
    \DisplayProof
    \\[2ex]
    \hline
  \end{array}\]
  \caption{قضیاتی رابطیاں لئی اگیتے والے !!{tableau} دے قاعدے}
  \ollabel{tab:prop-rules}
\end{table}

شاخ بند ہون دی شرط اوہی اے جیہڑی عام !!{tableau}s لئی اے،
پر !!{formula}s دے نال نال اگیتے وی اکو جیہے ہونے چاہیدے نیں۔
سو شاخ بند اے جے اوہدے وچ ایہ دونیں ہون:
\[
\sFmla{\True}{!A}[\sigma] \quad\text{تے}\quad \sFmla{\False}{!A}[\sigma]
\]
ایتھے $\sigma$ کوئی اگیتا اے تے~$!A$ کوئی !!{formula} اے۔

مفروضے رکھن دے قاعدے وی عام !!{tableau}s والے ای نیں، بس
مفروضیاں لئی اسیں ہمیشہ اگیتا~$1$ ورتدے آں۔ (اگیتا کوئی وی ہو سکدا اے،
جے سارے مفروضیاں دا اگیتا اکّو ہووے۔) سو مثلاً اسیں آکھدے آں کہ
\[
!B_1, \dots, !B_n \Proves !A
\]
اُدوں تے صرف اُدوں جدوں ہِیٹھلے مفروضیاں دا اک بند تابلو ہووے:
\[
\sFmla{\True}{!B_1}[1], \dots, \sFmla{\True}{!B_n}[1],
\sFmla{\False}{!A}[1].
\]

امکانی عمل کار\iftag{prvBox}{\iftag{prvDiamond}{~$\Box$
    تے}{~$\Box$}}{}\iftag{prvDiamond}{~$\Diamond$}{} لئی،
اگیتے~$\sigma$ والے !!a{formula} اُتے قاعدہ لاؤن نال جیہڑا نتیجہ
ملدا اے، اوہدا اگیتا $\sigma.n$ اے۔ پر کہڑا $n$ ورت سکدے آں،
ایہ ایس گل اُتے اے کہ نشان~$\True$ اے یا~$\False$۔

\iftag{prvBox}{قاعدہ $\TRule{\True}{\Box}$ اوہ شاخ ودھاؤندا اے
  جیہدے وچ $\sFmla{\True}{\Box !A}[\sigma]$ ہووے۔ ایہ اوہدے اُتے
  $\sFmla{\True}{!A}[\sigma.n]$ لاؤندا اے۔\iftag{prvDiamond}{ ایس طرح، }{}}{}\iftag{prvDiamond}{قاعدہ $\TRule{\False}{\Diamond}$ اوہ شاخ
  ودھاؤندا اے جیہدے وچ $\sFmla{\False}{\Diamond !A}[\sigma]$ ہووے۔
  ایہ اوہدے اُتے $\sFmla{\False}{!A}[\sigma.n]$ لاؤندا اے۔}{}
\iftag{notprvBox,notprvDiamond}{ایہ قاعدہ}{ایہ قاعدے} اوہدوں ای
لائے جا سکدے نیں جدوں اگیتا~$\sigma.n$ اوس شاخ اُتے \emph{پہلاں ای}
موجود ہووے جیہدے اُتے قاعدہ لگدا اے۔ شاخ لئی ایسے اگیتے نوں
اسیں «ورتیا ہویا» آکھدے آں۔

\iftag{prvBox}{قاعدہ $\TRule{\False}{\Box}$ اوہ شاخ ودھاؤندا اے
  جیہدے وچ $\sFmla{\False}{\Box !A}[\sigma]$ ہووے۔ ایہ اوہدے اُتے
  $\sFmla{\False}{!A}[\sigma.n]$ لاؤندا اے۔\iftag{prvDiamond}{ ایس طرح، }{}}{}\iftag{prvDiamond}{قاعدہ $\TRule{\True}{\Diamond}$ اوہ شاخ
  ودھاؤندا اے جیہدے وچ $\sFmla{\True}{\Diamond !A}[\sigma]$ ہووے۔
  ایہ اوہدے اُتے $\sFmla{\True}{!A}[\sigma.n]$ لاؤندا اے۔}{}
پر \iftag{notprvBox,notprvDiamond}{ایہ قاعدہ}{ایہ قاعدے} اوہدوں ای
لائے جا سکدے نیں جدوں اگیتا~$\sigma.n$ اوس شاخ اُتے پہلاں توں
\emph{موجود نہ} ہووے جیہدے اُتے قاعدہ لگدا اے۔ شاخ لئی ایسے اگیتے
نوں اسیں «نواں» آکھدے آں۔

قاعدے \olref{tab:rules-K} وچ دتے نیں۔

\begin{table}
  \begin{center}
    \def\arraystretch{3}\def\fCenter{}
    \begin{tabular}{|c|c|}
    \hline
    \iftag{prvBox}{
      \AxiomC{\sFmla{\True}{\Box !A}[\sigma]}
      \RightLabel{\TRule{\True}{\Box}}
      \UnaryInfC{\sFmla{\True}{!A}[\sigma.n]}
      \DisplayProof
      &
      \AxiomC{\sFmla{\False}{\Box !A}[\sigma]}
      \RightLabel{\TRule{\False}{\Box}}
      \UnaryInfC{\sFmla{\False}{!A}[\sigma.n]}
      \DisplayProof\\
      $\sigma.n$ ورتیا ہویا اے & $\sigma.n$ نواں اے
      \\[1ex]
      \hline}{}
    \iftag{prvDiamond}{
      \AxiomC{\sFmla{\True}{\Diamond !A}[\sigma]}
      \RightLabel{\TRule{\True}{\Diamond}}
      \UnaryInfC{\sFmla{\True}{!A}[\sigma.n]}
      \DisplayProof
      &
      \AxiomC{\sFmla{\False}{\Diamond !A}[\sigma]}
      \RightLabel{\TRule{\False}{\Diamond}}
      \UnaryInfC{\sFmla{\False}{!A}[\sigma.n]}
      \DisplayProof\\
      $\sigma.n$ نواں اے & $\sigma.n$ ورتیا ہویا اے
      \\[1ex]
      \hline}{}
    \end{tabular}
  \end{center}
  \caption{\Ax{K} دے امکانی قاعدے۔}
  \ollabel{tab:rules-K}
\end{table}

قاعدے \iftag{prvBox}{\TRule{\True}{\Box}}{\TRule{\False}{\Diamond}}
لئی شرط اے کہ اگیتا ورتیا ہویا ہووے۔ ایہ شرط ضروری اے، کیوں جے ایہدے
بغیر ہِیٹھلے !!{tableau} نوں بند گنیا جاندا:
\iftag{notprvBox,notprvDiamond}{%
  \iftag{prvBox}{%
  \begin{oltableau}
    [\pFmla{\True}{\Box \formula{A}}{1}, just = \TAss
      [\pFmla{\False}{\lnot\Box\lnot \formula{A}}{1}, just = \TAss
        [\pFmla{\True}{\formula{A}}{1.1}, just = {\TRule{\True}{\Box}[1]}
          [\pFmla{\True}{\Box\lnot\formula{A}}{1},
            just ={\TRule{\False}{\lnot}[2]}
            [\pFmla{\True}{\lnot\formula{A}}{1.1},
              just = {\TRule{\True}{\Box}[4]}
              [\pFmla{\False}{\formula{A}}{1.1},
                  just ={\TRule{\True}{\lnot}[5]}, close]
            ]
          ]
        ]
      ]
    ]
  \end{oltableau}
  }{
  \begin{oltableau}
    [\pFmla{\True}{\lnot\Diamond\lnot \formula{A}}{1}, just = \TAss
      [\pFmla{\False}{\Diamond\formula{A}}{1}, just = \TAss
        [\pFmla{\False}{\formula{A}}{1.1}, just = {\TRule{\False}{\Diamond}[2]}
          [\pFmla{\False}{\Diamond\lnot\formula{A}}{1},
            just ={\TRule{\True}{\lnot}[1]}
            [\pFmla{\False}{\lnot\formula{A}}{1.1},
              just = {\TRule{\False}{\Diamond}[4]}
              [\pFmla{\True}{\formula{A}}{1.1},
                  just ={\TRule{\True}{\lnot}[5]}, close]
            ]
          ]
        ]
      ]
    ]
  \end{oltableau}
}}{
  \begin{oltableau}
    [\pFmla{\True}{\Box \formula{A}}{1}, just = \TAss
      [\pFmla{\False}{\Diamond \formula{A}}{1}, just = \TAss
        [\pFmla{\True}{\formula{A}}{1.1}, just = {\TRule{\True}{\Box}[1]}
          [\pFmla{\False}{\formula{A}}{1.1},
            just = {\TRule{\False}{\Diamond}[2]}, close]
        ]
      ]
    ]
\end{oltableau}}

پر $\Box \formula{A} \Entails/ \Diamond \formula{A}$، سو ساڈا
اخذ دا نظام صحیح نہ رہندا۔ ایس طرح $\Diamond \formula{A} \Entails/
\Box \formula{A}$، پر جے قاعدے
\iftag{prvBox}{\TRule{\False}{\Box}}{\TRule{\True}{\Diamond}}
لئی اگیتا نواں ہون دی شرط نہ ہووے، تاں ایہ بند تابلو ہووے گا: \iftag{notprvBox,notprvDiamond}{%
  \iftag{prvBox}{%
  \begin{oltableau}
    [\pFmla{\True}{\lnot\Box\lnot \formula{A}}{1}, just = \TAss
      [\pFmla{\False}{\Box \formula{A}}{1}, just = \TAss
        [\pFmla{\False}{\formula{A}}{1.1}, just = {\TRule{\True}{\Box}[2]}
          [\pFmla{\False}{\Box\lnot\formula{A}}{1},
            just ={\TRule{\True}{\lnot}[1]}
            [\pFmla{\False}{\lnot\formula{A}}{1.1},
              just = {\TRule{\False}{\Box}[4]}
              [\pFmla{\True}{\formula{A}}{1.1},
                  just ={\TRule{\False}{\lnot}[5]}, close]
            ]
          ]
        ]
      ]
    ]
  \end{oltableau}
  }{
  \begin{oltableau}
    [\pFmla{\True}{\Diamond \formula{A}}{1}, just = \TAss
      [\pFmla{\False}{\lnot\Diamond\lnot\formula{A}}{1}, just = \TAss
        [\pFmla{\True}{\formula{A}}{1.1}, just = {\TRule{\True}{\Diamond}[1]}
          [\pFmla{\True}{\Diamond\lnot\formula{A}}{1},
            just ={\TRule{\False}{\lnot}[2]}
            [\pFmla{\True}{\lnot\formula{A}}{1.1},
              just = {\TRule{\True}{\Diamond}[4]}
              [\pFmla{\False}{\formula{A}}{1.1},
                  just ={\TRule{\True}{\lnot}[5]}, close]
            ]
          ]
        ]
      ]
    ]
  \end{oltableau}
}}{
  \begin{oltableau}
    [\pFmla{\True}{\Diamond \formula{A}}{1}, just = \TAss,name=one
      [\pFmla{\False}{\Box \formula{A}}{1}, just = \TAss, name=two
        [\pFmla{\True}{\formula{A}}{1.1}, just = {\TRule{\True}{\Diamond}[1]}
          [\pFmla{\False}{\formula{A}}{1.1},
            just = {\TRule{\False}{\Box}[2]}, close]
        ]
      ]
    ]
  \end{oltableau}}

\end{document}
